\documentclass[11pt,a4paper]{article}
\usepackage[margin=1in]{geometry}
\usepackage{amsmath,amssymb,amsthm}
\usepackage{algorithm}
\usepackage{algpseudocode}
\usepackage{booktabs}
\usepackage{hyperref}

\theoremstyle{definition}
\newtheorem{definition}{\textbf{Definition}}[section]
\newtheorem{theorem}{\textbf{Theorem}}[section]
\newtheorem{lemma}[theorem]{\textbf{Lemma}}
\newtheorem{remark}{\textbf{Remark}}[section]

\title{\textbf{Neural Radiance Fields and Volume Rendering}}
\author{\textbf{Team Scaibu} \\ \small Email: \texttt{jaydeep.9664920749@gmail.com} \\ \small \textbf{Architecture and Core Engine Team}}
\date{\textbf{October 2026}}

\begin{document}
\maketitle

\section{Definitions and Cognitive Mental Models}

\subsection{Formal Mathematical Definition}
\textbf{Definition (Continuous Radiance Field and Classical Volume Quadrature):}
Let a 3D scene be represented as a continuous function $F_\theta: (\mathbf{x}, \mathbf{d}) \to (\mathbf{c}, \sigma)$ parameterized by a multi-layer perceptron, where $\mathbf{x} = (x, y, z) \in \mathbb{R}^3$ is 3D spatial position, $\mathbf{d} = (\theta, \phi) \in \mathbb{S}^2$ is viewing direction, $\mathbf{c} = (r, g, b) \in [0, 1]^3$ is emitted radiance, and $\sigma \in \mathbb{R}_{\ge 0}$ is differential volume density.

High-frequency geometric detail is resolved by projecting coordinates through sinusoidal \textbf{Positional Encodings}:
{\boldmath
\begin{equation}
\gamma(p) = \left[ \sin(2^0 \pi p), \cos(2^0 \pi p), \, \dots, \, \sin(2^{L-1} \pi p), \cos(2^{L-1} \pi p) \right] \in \mathbb{R}^{2L}
\end{equation}
}
For a camera ray $\mathbf{r}(t) = \mathbf{o} + t \mathbf{d}$ with near and far bounds $[t_n, t_f]$, the expected rendered color is given by the continuous volume rendering integral:
{\boldmath
\begin{equation}
C(\mathbf{r}) = \int_{t_n}^{t_f} T(t) \, \sigma(\mathbf{r}(t)) \, \mathbf{c}(\mathbf{r}(t), \mathbf{d}) \, \mathrm{d}t, \quad T(t) = \exp\left( -\int_{t_n}^t \sigma(\mathbf{r}(s)) \mathrm{d}s \right)
\end{equation}
}
where $T(t)$ denotes cumulative transmittance from $t_n$ to $t$.

Under stratified discrete numerical quadrature across $N$ sample points with intervals $\delta_i = t_{i+1} - t_i$, the continuous integral simplifies to piecewise-constant alpha compositing:
{\boldmath
\begin{equation}
\hat{C}(\mathbf{r}) = \sum_{i=1}^N w_i \mathbf{c}_i, \quad w_i = T_i \left( 1 - \exp(-\sigma_i \delta_i) \right), \quad T_i = \prod_{j=1}^{i-1} \left( 1 - \alpha_j \right) = \exp\left( -\sum_{j=1}^{i-1} \sigma_j \delta_j \right)
\end{equation}
}
where $\alpha_i = 1 - \exp(-\sigma_i \delta_i)$ represents point opacity.

\subsection{Expanded Non-Technical Definition (Intuition for Anyone)}
\textbf{Intuitive Conceptual Definition:}
\begin{enumerate}
    \item \textbf{Replacing 3D Polygons with a Neural Fog}: Traditional computer graphics represent 3D worlds with millions of hard triangular meshes (polygons). NeRF replaces polygons with a continuous glowing cloud of colored fog stored inside an AI's brain.
    \item \textbf{Asking the AI at Any Coordinate}: If you ask the AI: ``What is at position $(x=1.2, y=3.4, z=5.6)$ if I look from the left?'', the AI answers instantly: ``There is red glowing matter with a density of 8.5.''
    \item \textbf{Marching Rays of Light}: To render a single pixel on your computer screen, you shoot an imaginary laser beam (ray) from the virtual camera through that pixel into the 3D scene.
    \item \textbf{Alpha Compositing (Adding Up the Fog)}: As the laser travels through the fog, it samples colors and densities along the way. Thick objects block the laser and add solid color; thin fog lets light pass through (transmittance). Combining all samples along the laser produces a stunning, photorealistic pixel with perfect shadows, reflections, and lighting!
\end{enumerate}

\subsection{Child-Friendly Mental Model (Physical Metaphor)}
Imagine looking through a series of stained-glass windows:
\begin{itemize}
    \item \textbf{Clear Glass}: Density is zero; all light passes straight through to the next layer.
    \item \textbf{Colored Translucent Glass}: Density is medium; it tints the light yellow and lets some pass through.
    \item \textbf{Solid Brick Wall}: Density is huge; it stops the light completely and reflects bright red brick.
    \item \textbf{The Composite View}: What your eye sees is the exact blend of all layers up to the brick wall!
\end{itemize}

\section{Core Mathematical Equations and Definitions}

\subsection{High-Frequency Positional Encoding Operator}
{\boldmath
\begin{equation}
\gamma(p) = \left[ \sin(2^0 \pi p), \cos(2^0 \pi p), \dots, \sin(2^{L-1} \pi p), \cos(2^{L-1} \pi p) \right]^T
\end{equation}
}
\begin{itemize}
    \item \textbf{Formal Definition}: The deterministic Fourier feature map expanding 1D spatial coordinates into $2L$ orthogonal harmonic bases.
    \item \textbf{What It Means}: Overcomes the spectral bias of deep MLPs, enabling standard fully-connected networks to learn high-frequency textures rather than oversmoothed blobs.
    \item \textbf{Why It Is Important}: Essential for rendering sharp object boundaries, hair, and complex reflections.
\end{itemize}

\subsection{Discrete Transmittance and Opacity Quadrature}
{\boldmath
\begin{equation}
\alpha_i = 1 - \exp(-\sigma_i \delta_i), \quad T_i = \prod_{j=1}^{i-1} (1 - \alpha_j) = \exp\left( -\sum_{j=1}^{i-1} \sigma_j \delta_j \right)
\end{equation}
}
\begin{itemize}
    \item \textbf{Formal Definition}: The discrete approximation of optical extinction along a ray segment of length $\delta_i$.
    \item \textbf{What It Means}: $\alpha_i$ is the probability that a photon is absorbed by segment $i$; $T_i$ is the probability that the photon survived all preceding segments without being absorbed.
    \item \textbf{Why It Is Important}: Guarantees that volume rendering adheres to physical Beer-Lambert optical attenuation.
\end{itemize}

\subsection{Discrete Compositing Weight Equation}
{\boldmath
\begin{equation}
w_i = T_i \alpha_i = \exp\left( -\sum_{j=1}^{i-1} \sigma_j \delta_j \right) \left( 1 - \exp(-\sigma_i \delta_i) \right)
\end{equation}
}
\begin{itemize}
    \item \textbf{Formal Definition}: The quadrature weight assigned to color sample $\mathbf{c}_i$ in the composited ray sum.
    \item \textbf{What It Means}: Peak weight occurs at the first dense surface encountered along the ray; occluded regions behind the surface receive near-zero weight.
    \item \textbf{Why It Is Important}: Enables precise surface depth estimation and novel-view occlusion handling.
\end{itemize}

\subsection{Expected Ray Depth Estimator}
{\boldmath
\begin{equation}
\hat{d}(\mathbf{r}) = \sum_{i=1}^N w_i t_i
\end{equation}
}
\begin{itemize}
    \item \textbf{Formal Definition}: The expectation of metric distance $t$ along the camera ray under the compositing weight distribution.
    \item \textbf{What It Means}: Produces a continuous, metric depth map of the scene directly from the radiance field.
    \item \textbf{Why It Is Important}: Allows extracting 3D point clouds and depth geometry for robotics and digital twin synthesis.
\end{itemize}

\section{Rigorous Mathematical Analysis Checklist}

\subsection{Notation and Symbol Semantics}
\begin{itemize}
    \item $\mathbf{r}(t) = \mathbf{o} + t \mathbf{d}$: Camera ray with origin $\mathbf{o}$ and direction $\mathbf{d}$.
    \item $\sigma_i \in \mathbb{R}_{\ge 0}$: Evaluated volume density.
    \item $\mathbf{c}_i \in [0, 1]^3$: Evaluated view-dependent RGB color.
    \item $\delta_i = t_{i+1} - t_i \in \mathbb{R}^+$: Step distance between quadrature points.
    \item $T_i \in [0, 1]$: Accumulated transmittance.
    \item $w_i \in [0, 1]$: Quadrature compositing weight.
\end{itemize}

\subsection{Epistemic Status Table}
\begin{table}[h]
\centering
\small
\begin{tabular}{@{}llll@{}}
\toprule
\textbf{Quantity} & \textbf{Symbol} & \textbf{Epistemic Status} & \textbf{Operational Role} \\
\midrule
Volume Density & $\sigma$ & Neural Observable & Quantifies local material opacity \\
Color Radiance & $\mathbf{c}$ & View-Dependent Neural Map & Emitted color under viewing direction $\mathbf{d}$ \\
Transmittance & $T_i$ & Exact Optical Integral & Tracks remaining light survival probability \\
Composited Color & $\hat{C}(\mathbf{r})$ & Differentiable Rendered Value & Compared directly against training photo pixels \\
\bottomrule
\end{tabular}
\end{table}

\subsection{Computational Dependency Graph}
{\centering
$\{\mathbf{o}, \mathbf{d}\} \longrightarrow \{t_i, \mathbf{x}_i\} \longrightarrow \text{MLP} \longrightarrow \{\sigma_i, \mathbf{c}_i\} \longrightarrow \{\alpha_i, T_i\} \longrightarrow w_i \longrightarrow \hat{C}(\mathbf{r}) \longrightarrow \mathcal{L}_{\text{photo}}$
\par}

\subsection{Dimension and Coordinate Consistency}
Spatial coordinates $\mathbf{x} \in \mathbb{R}^3$; ray parameters $t, \delta_i \in \mathbb{R}$; colors $\mathbf{c}_i, \hat{C} \in [0, 1]^3$; density $\sigma_i$ has units $[\text{distance}]^{-1}$.

\subsection{Asymptotic Regimes}
\begin{itemize}
    \item \textbf{Vacuum Limit ($\sigma \to 0$)}: $T_N \to 1, w_i \to 0$; total opacity vanishes, rendering empty transparent space.
    \item \textbf{Solid Surface Limit ($\sigma_i \to \infty$)}: $\alpha_i \to 1, T_{i+1} \to 0$; surface stops the ray completely, rendering opaque color.
\end{itemize}

\subsection{Boundary Constraints and Invariants}
\begin{itemize}
    \item \textbf{Partition of Unity}: $\sum_{i=1}^N w_i + T_{N+1} \equiv 1$ strictly holds for all density configurations.
\end{itemize}

\subsection{Structural Isomorphisms}
Volume rendering is structurally isomorphic to Continuous-Time Markov Jump Processes and survival analysis under hazard rate $\sigma(t)$.

\section{Theoretical Theorems and Formal Proofs}

\begin{theorem}[Exact Partition of Unity in Discrete Volume Rendering]
Let $\alpha_i = 1 - \exp(-\sigma_i \delta_i)$ with $\sigma_i \ge 0, \delta_i > 0$. Let $T_1 = 1$ and $T_{i+1} = T_i (1 - \alpha_i)$. Then the sum of quadrature weights plus the terminal transmittance equals unity identically:
{\boldmath
\begin{equation}
\sum_{i=1}^N w_i + T_{N+1} = 1
\end{equation}
}
\end{theorem}

\begin{proof}
Recall the definition of the quadrature weight:
{\boldmath
\begin{equation}
w_i = T_i \alpha_i = T_i (1 - (1 - \alpha_i)) = T_i - T_i(1 - \alpha_i)
\end{equation}
}
From the recurrence relation $T_{i+1} = T_i (1 - \alpha_i)$, we can rewrite $w_i$ as a telescoping difference:
{\boldmath
\begin{equation}
w_i = T_i - T_{i+1}
\end{equation}
}
Summing the weights from $i = 1$ to $N$:
{\boldmath
\begin{equation}
\sum_{i=1}^N w_i = \sum_{i=1}^N (T_i - T_{i+1}) = (T_1 - T_2) + (T_2 - T_3) + \dots + (T_N - T_{N+1})
\end{equation}
}
All intermediate terms cancel telescopically:
{\boldmath
\begin{equation}
\sum_{i=1}^N w_i = T_1 - T_{N+1}
\end{equation}
}
Since by definition $T_1 = 1$ (the ray enters the volume with zero initial attenuation):
{\boldmath
\begin{equation}
\sum_{i=1}^N w_i = 1 - T_{N+1} \iff \sum_{i=1}^N w_i + T_{N+1} = 1
\end{equation}
}
Thus, the composited weights form a valid defective probability distribution whose defect $T_{N+1}$ represents light passing unimpeded through to the background.
\end{proof}

\begin{theorem}[Equivalence of Discrete Alpha Compositing to Classical Over Operator]
The discrete volume quadrature recurrence $C_i = C_{i-1} + w_i \mathbf{c}_i$ is equivalent to Porter-Duff front-to-back alpha blending:
{\boldmath
\begin{equation}
\mathbf{C}_{\text{accum}} \gets \mathbf{C}_{\text{accum}} + T_{\text{accum}} \alpha_i \mathbf{c}_i, \quad T_{\text{accum}} \gets T_{\text{accum}}(1 - \alpha_i)
\end{equation}
}
\end{theorem}

\begin{proof}
Let $\mathbf{C}_0 = \mathbf{0}$ and $T_1 = 1$.
At step $i = 1$:
{\boldmath
\begin{equation}
\mathbf{C}_1 = \mathbf{0} + T_1 \alpha_1 \mathbf{c}_1 = \alpha_1 \mathbf{c}_1, \quad T_2 = 1 - \alpha_1
\end{equation}
}
At step $i = 2$:
{\boldmath
\begin{equation}
\mathbf{C}_2 = \mathbf{C}_1 + T_2 \alpha_2 \mathbf{c}_2 = \alpha_1 \mathbf{c}_1 + (1 - \alpha_1)\alpha_2 \mathbf{c}_2, \quad T_3 = (1 - \alpha_1)(1 - \alpha_2)
\end{equation}
}
By mathematical induction, at step $k$:
{\boldmath
\begin{equation}
\mathbf{C}_k = \sum_{i=1}^k \left( \prod_{j=1}^{i-1} (1 - \alpha_j) \right) \alpha_i \mathbf{c}_i = \sum_{i=1}^k T_i \alpha_i \mathbf{c}_i = \sum_{i=1}^k w_i \mathbf{c}_i
\end{equation}
}
This is identically the front-to-back optical over operator used in real-time GPU volume ray-marching.
\end{proof}

\section{Foundational Architectural Questions and Epistemic Meta-Questions}

\subsection{Architectural Question 1: NeRF vs 3D Gaussian Splatting}
\textbf{Question:} Why does 3D Gaussian Splatting render at 100+ FPS, while vanilla NeRF renders at 0.1 FPS?\\
\textbf{Answer:} NeRF queries a deep multi-layer perceptron 64 to 128 times along every single camera ray ($800 \times 800 \times 128 \approx 82 \text{ million}$ forward passes per image). 3D Gaussian Splatting replaces the neural network with explicit 3D Gaussian primitives, projecting them into screen space and sorting them into GPU tiles for instantaneous rasterization.

\textbf{Meta-Question:} Does explicit Gaussian splatting lose any representation capacity compared to NeRF?\\
\textbf{Answer:} Gaussians require substantial memory (millions of Gaussians can consume 1 to 2 GB of VRAM), whereas a NeRF MLP is a compact 5 MB neural file. NeRF represents a continuous neural field; Gaussian splatting is an explicit point-based volumetric proxy.

\subsection{Architectural Question 2: View-Dependent Effects and Geometric Overfitting}
\textbf{Question:} Why does NeRF inject viewing direction $\mathbf{d}$ only in the final RGB layers rather than the initial density layers?\\
\textbf{Answer:} Volume density $\sigma$ must depend strictly on 3D spatial position $\mathbf{x}$ to enforce physical multiview consistency (an object does not disappear when viewed from another angle). If density depended on $\mathbf{d}$, the network could overfit by rendering separate 2D billboard images for each camera angle without discovering real 3D geometry.

\textbf{Meta-Question:} How does NeRF capture specular highlights (glossiness)?\\
\textbf{Answer:} By conditioning the final color emission layer on $\mathbf{d}$. The density remains fixed, while the emitted color varies smoothly with camera angle, faithfully capturing reflection physics.

\section{Algorithmic Implementation Specification}

\begin{algorithm}[H]
\caption{Discrete Volume Rendering Quadrature Ray-Marcher}
\begin{algorithmic}[1]
\Require Densities $\boldsymbol{\sigma} \in \mathbb{R}^N$, colors $\mathbf{C}_{\text{in}} \in \mathbb{R}^{N \times 3}$, intervals $\boldsymbol{\delta} \in \mathbb{R}^N$, optional depths $\mathbf{t} \in \mathbb{R}^N$
\Ensure Composited RGB $\mathbf{C} \in \mathbb{R}^3$, terminal transmittance $T_{\text{term}}$, total opacity $\alpha_{\text{total}}$, weights $\mathbf{w} \in \mathbb{R}^N$
\State $T \gets 1.0, \quad \alpha_{\text{total}} \gets 0.0, \quad d_{\text{exp}} \gets 0.0$
\State Initialize $\mathbf{C} \gets [0.0, 0.0, 0.0]$, weights array $\mathbf{w}$ of length $N$
\For{$i = 0$ \textbf{to} $N - 1$}
    \State $\sigma_i \gets \max(0.0, \boldsymbol{\sigma}[i])$
    \State $\delta_i \gets \max(0.0, \boldsymbol{\delta}[i])$
    \State $\alpha_i \gets 1.0 - \exp(-\sigma_i \cdot \delta_i)$
    \State $w_i \gets T \cdot \alpha_i$
    \State $\mathbf{w}[i] \gets w_i$
    \State $\alpha_{\text{total}} \gets \alpha_{\text{total}} + w_i$
    \State $\mathbf{C}[0] \gets \mathbf{C}[0] + w_i \cdot \mathbf{C}_{\text{in}}[i][0]$
    \State $\mathbf{C}[1] \gets \mathbf{C}[1] + w_i \cdot \mathbf{C}_{\text{in}}[i][1]$
    \State $\mathbf{C}[2] \gets \mathbf{C}[2] + w_i \cdot \mathbf{C}_{\text{in}}[i][2]$
    \If{$\mathbf{t}$ is provided} \State $d_{\text{exp}} \gets d_{\text{exp}} + w_i \cdot \mathbf{t}[i]$ \EndIf
    \State $T \gets T \cdot (1.0 - \alpha_i)$
\EndFor
\State \Return $(\mathbf{C}, T, \alpha_{\text{total}}, \mathbf{w})$
\end{algorithmic}
\end{algorithm}

\section{Big-$\Theta$ Complexity Analysis}

\begin{itemize}
    \item \textbf{Ray Quadrature Time Complexity}: $\Theta(N)$ scalar operations per camera ray for $N$ sample points.
    \item \textbf{Full Image Rendering Complexity}: For an image of $H \times W$ pixels, total runtime is $\Theta(H \cdot W \cdot N \cdot \mathcal{C}_{\text{MLP}})$.
    \item \textbf{Space Complexity}: $\Theta(N)$ memory to store sample weights along the active ray.
    \item \textbf{Work \& Depth}: Work is $\mathcal{W} = \Theta(H \cdot W \cdot N)$; parallel depth across rays and sample steps is $\mathcal{D} = \Theta(\log N)$.
\end{itemize}

\section{Single Canonical Repository Reference}

The complete deterministic scalar implementation, verification suite, and unit test benchmarks are published at:
\begin{center}
\url{https://github.com/Chief-Strategist-J/llm-obs-infra}
\end{center}

\end{document}
